\section{An arbitrary-partner degree barrier for classical extraction layers}
\label{sec:classical-factor-degree-barrier}

\paragraph{Status and provenance.}
D-TD58, question Q2fn, is a new degree restriction for an
actual factorwise quotient, with its entire external partner retained.
The complete saved proof and invoked source passages have been read
by the author and pass; complete root and independent whole-proof reviews also pass.
The parent assigned
the remaining Lie-family gates after the alternating cutoff. The
author combines the accepted Larsen--Tiep estimate with the exact
normalizer-trace and quotient-cover arguments of D-TD02--03.
No independent-origin claim is made for those accepted ingredients.

\paragraph{Named covers and constants.}
Throughout, $d$ denotes the natural dimension, not a Lie rank, and
$q$ is the usual group parameter in
\[
 \operatorname{SL}_d(q),\quad \operatorname{SU}_d(q),\quad
 \operatorname{Sp}_d(q),\quad \Omega_d^\epsilon(q),\quad
 \operatorname{Spin}_d^\epsilon(q).
 \tag{CFD.1}
\]
The symplectic dimension is even; the unitary natural field is
$\mathbf F_{q^2}$. In the orthogonal cases the usual sign is omitted
for odd dimension, odd dimension is used only at odd $q$, and Spin
is used only at odd $q$. We require $d\ge120$ and the indicated
group to be quasisimple. A represented factor below must be an
actual central quotient of one of these named covers. An exceptional
cover is not included merely because its simple quotient has the
same Lie-type name.

Set
\[
 \sigma_0=10^{-29},\qquad C_0=2\cdot10^{30},\qquad
 B_0=2^{20},\qquad c_0=B_0^{-1}.
 \tag{CFD.2}
\]
For a finite-order operator, $\nu$ is its degree minus its largest
eigenspace dimension, allowing every eigenvalue target.

\begin{theorem}[D-TD58]
Let $C\to\mathbf P^1$ be an actual connected finite $J$-Galois
cover, with a faithful irreducible complex coefficient representation
on $V$, $e=\dim V$, and a generating product-one inertia tuple
$T_1,\ldots,T_r$ satisfying
\[
                  \sum_i\operatorname{rank}(T_i-I)=2e+6.
 \tag{CFD.3}
\]
Suppose $V=W\otimes U$ is preserved factorwise and let $P$ be the
actual finite projective image on $W$. Suppose $P$ contains and
normalizes the projective image of a finite perfect group
$N=E_1\cdots E_t$ acting faithfully and irreducibly on $W$, where
the $E_j$ are commuting quasisimple factors with central
intersections. The corresponding irreducible tensor factors $W_j$
are faithful $E_j$-modules of degrees $a_j$. Assume that each $E_j$
is a central quotient of a group in (CFD.1), with parameters
$d_j\ge120,q_j$. Then the simultaneous inequalities
\[
                         a_j\ge q_j^{C_0d_j}
                         \quad(1\le j\le t)
 \tag{CFD.4}
\]
are impossible.

In particular, consider any one nonabelian coarse layer in a complete
D-TD12 decomposition of a finite primitive source satisfying
(CFD.3). If its isomorphic quasisimple components admit one of the
specified named-cover presentations with $d\ge120$, then every
component degree satisfies
\[
                              a<q^{C_0d}.
 \tag{CFD.5}
\]
The whole external partner can have arbitrary dimension and image.
No twelve-mark condition is imposed on the projected quotient,
and the conclusion does not require an alternating cutoff, equality
of coefficient fields, or a self-duality hypothesis.
\end{theorem}

\subsection*{1. Exact quantitative inputs and central quotients}

Larsen--Tiep's all-element estimate, with the smaller explicit
constant checked in D-UNW-01 and D-OL-01, gives for every irreducible
ordinary character $\chi$ of a named cover in (CFD.1)
\[
 \frac{|\chi(g)|}{\chi(1)}
       \le\chi(1)^{-\sigma_0\operatorname{supp}(g)/d}.
 \tag{CFD.6}
\]
Here support is $d$ minus the largest natural eigenspace dimension
over the algebraic closure; in the Spin case it is computed on the
natural orthogonal image. The exact source is
\cite[Theorem~5.5 and complete proof, accepted-author pp.~25--26]
{UNW:LT2024}. The supplied constant is
$[6507\cdot2^{21}\cdot10^{18}]^{-1}>7\cdot10^{-29}$.
Theorem~4.4 and its proof, author pp.~19--20, supply the constants
$2^{14}10^{12}$ and $(2^{13}10^{12})^{-1}$ used there.

We also use the ordinary minimum-degree consequence of the
Landazuri--Seitz bounds as explicitly invoked in that paper's proof
of Corollary~5.6, author p.~26. A deliberately weaker common
consequence, in the natural-dimension conventions (CFD.1), is
\[
 \theta(1)\ge q^{d/6}\ge B_0
 \quad\text{for every nontrivial ordinary irreducible character
 of the simple quotient, when }d\ge120.
 \tag{CFD.7}
\]
To check the convention, the source uses a bound $q^{n/3}$ for
$\operatorname{SL}_n(q),\operatorname{SU}_n(q),\Omega_n^\epsilon(q)$
and $\operatorname{Spin}_n^\epsilon(q)$, and lists the symplectic
group as $\operatorname{Sp}_{2n}(q)$. Taking $d=2n$ in that case
gives the weaker exponent $d/6$. The natural-field notation for
unitary groups is specified on author p.~4; using the ordinary
group parameter $q$ only weakens the stated floor. All parameters
here are outside the small-dimensional exceptions. Inflating a
simple-quotient character to its named cover preserves its degree.

The irreducible $E_j$-character also inflates irreducibly to its
named cover, with the same degree $a_j$. A noncentral element of
$E_j$ has only noncentral lifts. In a linear, unitary, symplectic or
natural orthogonal group, natural support zero means the element
is scalar and hence central. In Spin, a scalar natural image has
commutators in the central kernel of the natural map; these give a
homomorphism from the perfect group into that kernel and are
therefore trivial. Thus support is at least one on every relevant
noncentral lift.

Assuming (CFD.4), it follows that
\[
 \frac{|\chi_j(x)|}{a_j}
       \le q_j^{-\sigma_0 C_0}
       =q_j^{-20}\le c_0
       \qquad(x\in E_j\setminus Z(E_j)).
 \tag{CFD.8}
\]
The character of $N$ on $W$ is the tensor product of these
characters, with the central identifications respected. Hence it
has normalized absolute value at most $c_0$ on every noncentral
element of $N$. Put $A=\dim W$. The central quotient
$S=N/Z(N)$ is the direct product of the $E_j/Z(E_j)$, so every
nontrivial ordinary irreducible $S$-module has degree at least $B_0$.

\subsection*{2. Every outer or permuting normalizer}

Because $N$ is perfect, its represented determinant is one. The
finite special-linear inverse image of $P$ therefore contains $N$.
Moreover projective normalization is actual normalization: an
operator normalizing $\mathbf C^*N$ normalizes its commutator
subgroup $(\mathbf C^*N)'=N$. The special-linear inverse image is
finite with scalar kernel $\mu_A$, so all its operators and $N$
can be made unitary simultaneously.

Let $T$ be a nonscalar representative of an element of $P$, and
let $\alpha$ be its induced automorphism of $S$. If $\alpha=1$,
then $TxT^{-1}x^{-1}$ is central for every $x\in N$. The resulting
central discrepancy is a homomorphism on the perfect group $N$,
so it is trivial. Schur's lemma then makes $T$ scalar, a
contradiction. Thus $H=C_S(\alpha)$ is a proper subgroup.

The nonconstant part of the permutation module on $S/H$ contains
a nontrivial ordinary irreducible constituent. Equation (CFD.7)
therefore implies
\[
                  [S:H]\ge B_0+1,
                  \qquad |H|/|S|\le(B_0+1)^{-1}.
 \tag{CFD.9}
\]
Schur averaging on the irreducible $N$-module $W$ gives
\[
 \frac{|\operatorname{tr}T|^2}{A^2}
  =\frac1{|N|}\sum_{x\in N}
       \frac{\operatorname{tr}(TxT^{-1}x^{-1})}{A}
  \le c_0+\frac{1-c_0}{B_0+1}
  =\frac2{B_0+1}<\frac1{170^2}.
 \tag{CFD.10}
\]
For clarity, the exact averaging identity follows by averaging
$xT^{-1}x^{-1}$, which equals $(\operatorname{tr}T^{-1}/A)I$.
The commutator is central precisely on the inverse image of $H$.
Its normalized trace has real part at most one there, and absolute
value at most $c_0$ elsewhere, by (CFD.8). Taking real parts proves
the inequality even when a central commutator is a nonidentity
scalar. The last numerical comparison is the integer inequality
$2\cdot170^2=57800<2^{20}+1$.

This is the full-normalizer mechanism of D-TD03 written for the
new degree hypothesis. It includes outer automorphisms and
permutations of the quasisimple factors, without choosing separate
quotient homomorphisms to those permuted components.

\subsection*{3. The actual quotient and complete original periods}

Form the actual connected $P$-cover
$Z=C/\ker(J\to P)\to\mathbf P^1$. Its group contains $S$.
In particular it is nonsolvable and has order greater than 60:
a nontrivial ordinary character of a simple factor has degree at
least $B_0$, so character orthogonality already gives that factor
order greater than $B_0^2$. The genus argument of D-TD02 therefore
gives $g(Z)\ge2$. Namely a genus-zero finite automorphism group is
solvable or has order at most 60, while an elliptic automorphism
group is an extension of an abelian translation group by a cyclic
group. These are the exact already proved actual-cover interfaces.

Let $y_i$ be the image of the original inertia $T_i$ in $P$ and
let $o_i$ be its order when $y_i\ne1$. If the original coefficient
order is $L_i$, then $o_i\mid L_i$. Normalize the two tensor
representatives reciprocally to have orders dividing $L_i$: their
$L_i$-th powers are inverse scalars, so a choice of an $L_i$-th root
makes both powers the identity without changing $T_i$. Make them
unitary. Among the full $L_i$ powers, precisely $L_i/o_i$ have
scalar $W$-action. Every other power has normalized absolute
$W$-trace less than $1/170$ by (CFD.10); the entire $U$-trace
has normalized absolute value at most one. Fourier averaging at
the target eigenvalue one gives
\[
 \frac{\operatorname{rank}(T_i-I)}e
       \ge(1-1/170)(1-1/o_i)\qquad(y_i\ne1).
 \tag{CFD.11}
\]
No projected-period shortcut, scalar-order restriction, or condition
on the whole external partner has been inserted. Projected identity
entries keep their nonnegative actual costs and may be discarded.

The actual quotient Riemann--Hurwitz formula has positive orbifold
excess. The elementary integer-order bound proved in D-TD02 is
\[
             \sum_{y_i\ne1}(1-1/o_i)\ge2+1/42=85/42.
 \tag{CFD.12}
\]
Summing (CFD.11) and using (CFD.3) gives
\[
          2+6/e\ge(1-1/170)85/42,
          \qquad e(1-85/170)\le252,
          \qquad e\le504.
 \tag{CFD.13}
\]
But $e\ge A\ge a_j\ge q_j^{C_0d_j}>1000$. This contradiction
proves the first assertion.

\subsection*{4. Every represented layer in the complete recursion}

At a nonabelian minimal-normal step of D-TD12, the represented
perfect group on the extracted coarse factor is the central product
of its commuting quasisimple components. It is faithful and
irreducible there. Its projective image is contained and normalized
in an actual finite projective quotient of the original $J$.
Every original operator has a factorwise tensor expression with
the complete complementary factor, including all earlier and later
extractions. These are properties of the represented quotient
image; the components are not asserted to lift to normal subgroups
of the original $J$.

The components of this minimal-normal layer form one transitive
orbit, so their degrees are equal. Transport a chosen named-cover
presentation of one component through those actual isomorphisms.
This supplies the same $d,q$ presentation and the same degree $a$
for each component; no classification-based uniqueness of parameters
or invariance of a character level is assumed. If (CFD.5) failed,
all component inequalities (CFD.4) would hold and the preceding
argument would exclude the actual source. This proves the
recursive conclusion.

\paragraph{Strict scope and remaining boundary.}
The earlier D-OL-02 degree barrier required its selected large-degree
factors to account for the complete coefficient representation in
order to charge twelve original entries. The present argument permits
an arbitrary additional factor, through its actual quotient curve
and every original tensor period. It is not an assertion that the
previous numerical constant has improved within that earlier scope.
The new barrier bounds the logarithmic degree divided by natural
dimension on each indicated classical extraction layer. It does
not bound natural rank at small fields, classify the remaining
low-degree characters, cover unnamed exceptional central covers,
or yield a finite catalogue of all original sources.
